\documentclass[10pt,a4paper,ssfamily]{exam}
\usepackage[brazil]{babel}
\usepackage[numbers,sort&compress]{natbib}
\usepackage[utf8]{inputenc}
\renewcommand{\baselinestretch}{1.0}
\usepackage{epsfig}
\usepackage{mhchem}
\usepackage{graphicx}
\usepackage{makeidx}
\usepackage{multicol}
\usepackage{multirow}
\usepackage{amssymb}
\usepackage{amsmath}
\DeclareMathOperator{\sen}{sen}
\newcommand{\listfont}{\bf\fontencoding{OT1}\sffamily\large}
\newcommand{\titlesfont}{\fontencoding{OT1}\sffamily\large}
\renewcommand{\baselinestretch}{1.0}
\newcommand{\Mg}{Mg$^{2+}$}
\renewcommand{\t}{\textrm}
\newcommand{\cc}[1]{[\textrm{#1}]}
\usepackage{enumitem}
\newcommand{\bmat}{\left[\begin{array}{cc}}
\newcommand{\emat}{\end{array}\right]}
\newcommand{\nomera}{
  \noindent
  \begin{tabular}{|p{0.7\textwidth}|p{0.3\textwidth}|}
  Nome: & RA: \\
  \hline
  \end{tabular}
}

\newcommand{\qini}{\begin{enumerate}[resume]\item}
\newcommand{\qend}{\end{enumerate}}

\newcommand{\oC}{$^\circ$C}
\newcommand{\kcalmol}{~kcal~mol$^{-1}$}
\newcommand{\moll}{~mol~L$^{-1}$}
\newcommand{\molkg}{~mol~kg$^{-1}$}
\newcommand{\gl}{~g~L$^{-1}$}
\newcommand{\e}[1]{$\times 10^{#1}$}

\newcommand{\eps}{\varepsilon}

\newcommand{\Resposta}[1]{\vspace{0.5cm}\noindent{\bf Resposta:}\\ #1
\begin{center}\rule{\textwidth}{0.4pt}\end{center}}
\renewcommand{\Resposta}[1]{}
\newcommand{\resposta}[1]{\vspace{0.5cm}\noindent{\bf Resposta:}\\ #1
\begin{center}\rule{\textwidth}{0.4pt}\end{center}}


\begin{document}
\sffamily
\pagestyle{empty}
\nomera

\begin{center}

{\bf QG101 - Química Geral}

{\bf Aula 1 - Medidas, Propriedades e o Mol}\\
\underline{Justifique suas respostas.}
\end{center}

\begin{center}{\bf Questões}\end{center}

% ---------------------------------------------------------------
\qini
Um engenheiro ambiental analisa a concentração de chumbo (\ce{Pb}) em
amostras de água. O laudo indica $15{,}0~\mu\text{g/L}$ de \ce{Pb}.
\begin{enumerate}
  \item Converta esse valor para mg/L e para g/m$^3$, mostrando os
        fatores de conversão utilizados.
  \item O limite máximo permitido pela legislação brasileira (CONAMA) é
        de $0{,}010~\text{mg/L}$. A amostra está dentro do limite?
        Justifique numericamente.
  \item A massa molar do \ce{Pb} é $207{,}2~\text{g/mol}$. Quantos
        \emph{nanomols} de \ce{Pb} existem em 1,00 L dessa amostra?
\end{enumerate}

\vspace{2cm}
\qend

% ---------------------------------------------------------------
\qini
Classifique cada grandeza abaixo como \textbf{extensiva (E)} ou
\textbf{intensiva (I)} e justifique sua classificação em uma frase:
\begin{enumerate}
  \item Temperatura de fusão do ferro a 1 atm.
  \item Número de mols de \ce{CO2} liberados por um motor a combustão
        durante 1 hora de operação.
  \item Densidade do etanol a 25\oC.
  \item Energia total armazenada em um tanque de combustível.
  \item Concentração molar de \ce{NaCl} em uma solução salina.
\end{enumerate}
\noindent
\textbf{Bônus:} dê um exemplo de como duas grandezas extensivas da lista
acima poderiam ser combinadas para gerar uma grandeza intensiva.

\vspace{2cm}
\qend

% ---------------------------------------------------------------
\qini
A glicose (\ce{C6H12O6}, $M = 180{,}16~\text{g/mol}$) é o principal
combustível celular do corpo humano.
\begin{enumerate}
  \item Quantos mols de glicose existem em $270{,}2~\text{g}$?
  \item Quantas \emph{moléculas} de glicose isso representa?
        Expresse sua resposta em notação científica.
  \item A glicose contém 6 átomos de carbono por molécula.
        Quantos átomos de carbono existem nessa amostra?
  \item Se essa quantidade de glicose fosse dissolvida em água até
        completar $500~\text{mL}$ de solução, qual seria a concentração
        molar da solução?
\end{enumerate}

\vspace{2cm}
\qend

\end{document}
